\documentclass[sigconf,review,anonymous]{acmart}

\sloppy
\emergencystretch=1em
\settopmatter{printacmref=true}
\renewcommand\footnotetextcopyrightpermission[1]{}
\pagestyle{plain}

\acmConference[SECUTE 2026]{2nd International Workshop on Software Security Testing}{October 16, 2026}{Munich, Germany}
\acmYear{2026}
\acmISBN{}
\acmDOI{}

\newcommand{\tool}{PolicyStrata}
\newcommand{\ArtifactSha}{d8182fb6d9b6d6a4a8ffec119e53ec6b3d34896bc4c7c24d38fb3ba376e76225}
\newcommand{\ArtifactUrl}{Anonymized artifact URL to be inserted after external anonymous hosting is created.}

\title{Finding Cross-Layer Authorization Drift in LLM Data-Agent Systems}

\author{Anonymous Author(s)}
\affiliation{%
  \institution{Anonymous Institution}
  \country{}}

\begin{document}

\begin{abstract}
LLM data agents split security policy across model-visible tools, constrained
generation grammars, semantic validators, SQL compilers, database roles, row-level
security, lineage, and release filters. A layer can satisfy its local checks
while the composed system exposes a forbidden metric, drops a tenant predicate,
or releases a result whose lineage should remain inside a trusted boundary. We
present \tool{}, a deterministic security-testing artifact for cross-layer policy
drift in governed data-agent stacks. Given canonical authorization, semantic,
and release specifications, \tool{} generates or imports principals, requests,
semantic plans, mixed-version surface configurations, SQL traces, database
states, and result-lineage decisions. It then checks responsibility-scoped
contracts and emits minimized witnesses for the first violated transition. The
artifact includes deterministic benchmark suites, baseline comparators, a stack
doctor for policy/TOS/privacy-document and prompt-manifest coverage, and an
anonymous security fixture covering unauthorized capability exposure, stale
tenant-key lowering, and release leakage after database containment.
\end{abstract}

\keywords{security testing, LLM data agents, authorization, text-to-SQL, row-level security, policy testing}

\begin{CCSXML}
<ccs2012>
   <concept>
       <concept_id>10002978.10003006.10003013</concept_id>
       <concept_desc>Security and privacy~Software security engineering</concept_desc>
       <concept_significance>500</concept_significance>
   </concept>
   <concept>
       <concept_id>10011007.10011074.10011099</concept_id>
       <concept_desc>Software and its engineering~Software testing and debugging</concept_desc>
       <concept_significance>300</concept_significance>
   </concept>
</ccs2012>
\end{CCSXML}

\ccsdesc[500]{Security and privacy~Software security engineering}
\ccsdesc[300]{Software and its engineering~Software testing and debugging}

\maketitle

\section{Problem}

LLM data agents translate user requests into tool calls, semantic plans, SQL
queries, and final results. Security-critical obligations are therefore repeated
across heterogeneous surfaces: the model-visible manifest says which operations
exist; a grammar or schema constrains action shape; a validator decides whether a
semantic object is authorized; a compiler lowers the object to executable SQL; a
database role and row-level security (RLS) policy contain execution; and a
release filter decides whether a result-lineage pair may leave the trusted
boundary.

This decomposition creates a software-security testing problem. A conventional
unit test can verify that a validator rejects a forbidden request, a SQL parser
accepts generated SQL, or RLS prevents cross-tenant reads. None of those checks
alone establishes that the same tenant, purpose, policy version, semantic
meaning, and release obligation survived every transition. Constrained decoding
systems such as PICARD reduce malformed actions~\cite{picard}, and text-to-SQL
benchmarks such as Spider and BIRD evaluate query correctness~\cite{spider,bird},
but grammar membership and query correctness are not authorization guarantees.
Likewise, PostgreSQL RLS is an important containment layer, but it does not
decide whether a result computed from sensitive lineage is safe to release
\cite{postgresql_rls}.

\tool{} targets this gap. It treats cross-layer authorization drift as a
security test case: a principal, request, policy version vector, semantic plan,
lowered SQL, database state, lineage, observed result, and violated obligation.
The goal is not to replace authorization enforcement. It is to generate and
replay tests that expose when deployed policy-bearing representations no longer
agree on their assigned obligations.

\section{Threat Model and Contracts}

We model accidental and adversarial drift in a data-agent stack. The attacker or
fault can influence a user request, exploit a stale model-visible tool manifest,
trigger a stale validator or compiler path, or route a result through a release
surface that lacks sufficient lineage checks. The trusted input is a canonical
policy snapshot. We do not assume the language model is trustworthy, and we do
not treat constrained generation as an enforcement boundary. We also do not
assume that database containment is enough: an unauthorized plan may be contained
by RLS while still revealing unsafe aggregate or lineage information through a
later release path.

\tool{} separates policy into authorization, semantic, and release
specifications. Each surface gets a responsibility-scoped contract:

\begin{itemize}
  \item \textbf{Manifest/grammar:} do not expose capabilities whose reachable
  operations are unauthorized for the principal and context.
  \item \textbf{Validator:} accept authorized semantic objects inside the
  declared support envelope and reject unauthorized objects.
  \item \textbf{Compiler:} preserve principal, tenant, purpose, policy version,
  lineage, and business semantics during lowering.
  \item \textbf{Database policy:} contain unauthorized executable operations
  under the runtime role.
  \item \textbf{Release policy:} release only result-lineage pairs allowed to
  leave the trusted boundary.
\end{itemize}

This contract split is similar in spirit to access-control policy mutation and
difference testing~\cite{xacml_mutation,margrave}, but the tested object is the
relationship among policy-bearing representations rather than a single policy
language. The policy oracle is independent of the SQL compiler path so a compiler
bug cannot certify itself.

\section{Artifact Workflow}

The artifact has two reviewer-facing workflows.

\textbf{Scanner.} The scanner imports JSONL tool/SQL traces or deterministic
benchmark cases, authorizes the semantic request against the canonical policy,
checks whether SQL preserves tenant scope and metric semantics, optionally
executes read-only SQL in a PostgreSQL fixture, checks RLS and database state
assertions, and evaluates release decisions. Findings are written as
\texttt{scan.json}, \texttt{findings.jsonl}, \texttt{summary.json},
\texttt{report.md}, and minimized witness files. Each witness records the
principal, request, version vector, semantic plan, lowered SQL, database rows or
state, lineage, observed result, and first violated transition.

\textbf{Doctor.} The doctor audits stack wiring before a team has complete test
coverage. It inventories policy/domain YAML, surface contracts, dbt semantic
models, SQL/tool traces, tenant predicates, PostgreSQL schema/RLS/grants/views,
state assertions, release tests, CI gates, privacy policies, terms of service,
data-processing agreements, security policies, internal access policies, prompt
manifests, and source maps. For privacy/TOS inputs, it classifies document type
and extracts deterministic obligation signals such as personal-data
minimization, purpose-limited processing, retention/deletion, third-party
sharing, subprocessor controls, security controls, sensitive-data controls, and
tenant isolation. For prompt manifests, it compares model-visible metrics,
dimensions, and tools against the canonical policy. The output is coverage
accounting plus remediation todos with owner, files, expected tests, and CI gate
command.

\section{Security Fixture}

The SECUTE artifact includes a focused security fixture under
\texttt{examples/secute\_security\_cases}. The fixture is intentionally small
enough to inspect by hand, but it exercises the surfaces that are easy to miss in
ordinary application tests: privacy/TOS/DPA/internal-policy inputs,
model-visible tools, source maps, semantic models, SQL traces, release
decisions, and tenant-scope lowering. Table~\ref{tab:security-cases} summarizes
the cases.

\begin{table}[t]
\caption{Anonymous SECUTE fixture cases.}
\label{tab:security-cases}
\begin{tabular}{p{0.25\columnwidth}p{0.38\columnwidth}p{0.25\columnwidth}}
\toprule
Case & Drift & Expected finding \\
\midrule
Unauthorized exposure & Manifest exposes a forbidden revenue alias that reaches
SQL and is released. & Unsafe release; unauthorized SQL reachability. \\
Stale tenant key & Compiler lowers an authorized metric through
\texttt{legacy\_tenant\_id}. & Tenant-scope lowering violation. \\
Release leak & Export releases customer-email lineage after DB containment
would have scoped rows. & Unsafe release; unauthorized SQL reachability. \\
\bottomrule
\end{tabular}
\end{table}

The fixture also checks governance wiring. The privacy policy and terms fixture
trigger personal-data minimization, purpose limitation, notice/consent,
retention/deletion, third-party sharing, security-control, sensitive-data, and
tenant-isolation obligations. The DPA fixture adds controller/processor and
subprocessor controls. The internal policy fixture adds least-privilege and
role-based-access signals. The prompt manifest intentionally exposes unsafe
metrics and sensitive dimensions, so doctor reports it as partial rather than
clean. The source map connects each failing trace to a notional owner and test
file, which lets the report produce remediation work instead of only listing
violations.

\section{Evaluation Snapshot}

The deterministic benchmark path contains three synthetic analytics domains and
four suite families: seeded cases, generated cases, detector-frozen held-out
cases, and clean controls. In one artifact reproduction run, \tool{} killed
1720/1720 non-clean injected faults covered by the implemented operators and
reported 0 false positives on 80 clean controls. The run produced 1800 traces and
the artifact archive is 211775 bytes. These are deterministic artifact-suite
coverage numbers, not production recall and not an authorization-boundary claim.

\begin{table}[t]
\caption{Point baselines over the same deterministic non-clean traces.}
\label{tab:baselines}
\begin{tabular}{lr}
\toprule
Comparator & Caught \\
\midrule
Grammar only & 121/1720 \\
Semantic validator only & 573/1720 \\
SQL policy checker & 1043/1720 \\
Database policy only & 326/1720 \\
Release filter only & 364/1720 \\
Defense-in-depth stack & 1561/1720 \\
\bottomrule
\end{tabular}
\end{table}

For SECUTE, we added an anonymous security fixture that exercises three concrete
failure classes. First, a model-visible revenue tool exposes a forbidden alias
that reaches SQL and is released. Second, a compiler lowers an authorized ticket
metric through a stale tenant key rather than the canonical tenant predicate.
Third, an export path releases customer-email lineage even though the request
should be withheld at the release boundary. Running the scanner on this fixture
produces five high-confidence findings: two unsafe releases, two unauthorized
requests reaching SQL, and one stale tenant-scope lowering. Running doctor on
the same fixture reports wired privacy/TOS/DPA/internal-policy ingestion, wired
source mapping, wired export coverage, and partial prompt-manifest checks because
the manifest intentionally exposes unsafe capabilities.

\section{Reproducibility}

The anonymous artifact is built as a zip with a generated manifest and an
anonymization scan. The builder omits public project URLs, public release notes,
GitHub Action branding, generated caches, local run outputs, and private local
paths. It generates an anonymous README, license, data-availability note, and
project metadata while preserving the package name used by imports and command
examples.

On a fresh unpack of the current archive, the following checks passed:
\texttt{uv run --extra dev ruff check .}; \texttt{uv run --extra dev mypy src};
and \texttt{uv run --extra dev pytest}. The test run collected 105 tests, with
100 passing and 5 skipped optional integration tests for PostgreSQL and
ClickHouse smoke paths. Running the SECUTE scan exits nonzero by design because
the fixture is a failing security-test suite. The expected gate output is:

\begin{scriptsize}
\begin{verbatim}
unsafe_release_unauthorized_metric_exposure_reached_sql
unauthorized_trace_reached_sql_unauthorized_metric_exposure_reached_sql
tenant_scope_missing_stale_tenant_key_lowering
unsafe_release_release_leak_after_database_containment
unauthorized_trace_reached_sql_release_leak_after_database_containment
\end{verbatim}
\end{scriptsize}

The main benchmark reproduction path is longer because it regenerates frozen
benchmark manifests, runs deterministic suites, computes baselines and
ablations, and writes evidence reports. Its purpose is reviewer verification,
not runtime benchmarking.

\section{Limitations and Ethics}

\tool{} is a testing artifact, not an access-control engine. It does not prove
general SQL equivalence, does not claim recall on unknown production incidents,
and does not replace runtime enforcement. The deterministic benchmark operators
can miss organization-specific policy nuance, and optional database fixtures are
smoke evidence unless backed by representative schema and state. The anonymous
security fixture uses synthetic data and synthetic policy documents. No customer
data, private prompts, or production secrets are included.

\section{Relation to Prior Work}

\tool{} composes ideas from several testing lines. Access-control mutation
testing supplies the adequacy intuition: a test suite should kill meaningful
policy mutants, while equivalent mutants must be separated from missed faults
\cite{xacml_mutation}. Policy-difference analysis shows how small policy
version changes can alter authorization behavior~\cite{margrave}. Text-to-SQL
test-suite accuracy shows why generated states and semantic comparison matter
for query behavior~\cite{test_suite_accuracy}. Database access control and RLS
remain important enforcement layers~\cite{postgresql_rls}. \tool{} uses these
ideas across the data-agent pipeline: the artifact is successful when it
identifies which representation first stopped carrying its assigned obligation.

\section{Data Availability Statement}

All data in the submitted artifact is synthetic or generated by deterministic
fixtures. The anonymous artifact archive contains source, tests, benchmark
generators, JSONL traces, optional PostgreSQL fixture files, expected outputs,
and the SECUTE security fixture. It requires Python 3.10 or newer and
\texttt{uv}; it does not require an LLM API key or host \texttt{psql}. Docker is
used only for optional database smoke checks.

Artifact access during review: \ArtifactUrl{} The local review archive checksum
is:

\begin{scriptsize}
\begin{verbatim}
SHA256 d8182fb6d9b6d6a4a8ffec119e53ec6b3d34896bc4c7c24d38fb3ba376e76225
\end{verbatim}
\end{scriptsize}

Primary commands:

\begin{scriptsize}
\begin{verbatim}
uv run --extra dev ruff check .
uv run --extra dev mypy src
uv run --extra dev pytest
POLICYSTRATA_RUN_ROOT=/tmp/policystrata-final \
  ./scripts/reproduce-final.sh
uv run policystrata scan --config \
  examples/secute_security_cases/policystrata.yaml \
  --out /tmp/policystrata-secute-scan
\end{verbatim}
\end{scriptsize}

\bibliographystyle{ACM-Reference-Format}
\bibliography{refs}

\end{document}
